\section{Discussion}

\subsection{What is currently convincing}

The most convincing current conclusion is that reconstruction-based anomaly detection is viable for this dataset and that \texttt{cons\_hostatgeria\_underfloor\_hea} provides a strong thesis case. It is physically distinctive, it overlaps the engineering baseline, and its anomaly patterns remain interpretable under both raw and stabilized flow views.

\subsection{What remains uncertain}

The biggest remaining limitation is still validation. Most sheets do not have labeled faults, so anomaly interpretation depends on triangulation:

\begin{itemize}
\item reconstruction error
\item dominant feature
\item signal plots
\item low delta-T overlap
\item supervisor and domain review
\end{itemize}

This means the thesis should frame most anomalies as credible candidates rather than confirmed faults.

\subsection{Implications for live deployment}

The historical result pack is already useful for deployment preparation. Once one to two months of live data become available, the intended strategy is to apply the historical-trained detector prospectively and evaluate the flagged windows through operational review rather than standard classification metrics.

\subsection{Writing priorities}

When drafting the final thesis, the narrative should remain focused. The autoencoder is the main method. The low delta-T baseline is the main engineering comparison. Dominant-feature attribution is the main anomaly interpretation layer. Clustering should be presented as a supporting analysis rather than the core detection method.
